\documentclass[conference,a4paper]{IEEEtran}
\IEEEoverridecommandlockouts
\usepackage{cite}
\usepackage{amsmath,amssymb}
\usepackage{graphicx}
\usepackage[hidelinks]{hyperref}

\begin{document}

\title{Administrative Price Ceilings and the Merchant Storage Business Case:
Evidence from India's Day-Ahead Power Exchange}

\author{
\IEEEauthorblockN{AVR Akshaya}
\IEEEauthorblockA{\textit{School of Business and Management}\\
\textit{CHRIST (Deemed to be University)}\\
Pune, India\\
akshaya.avr@christuniversity.in}
\and
\IEEEauthorblockN{Jerry Silvon J M}
\IEEEauthorblockA{\textit{School of Business and Management}\\
\textit{CHRIST (Deemed to be University)}\\
Pune, India\\
jerry.silvon@bbabah.christuniversity.in}
}

\maketitle

\begin{abstract}
India's Central Electricity Regulatory Commission cut the day-ahead price
ceiling from Rs20/kWh to Rs12/kWh in April 2022 and to Rs10/kWh in April 2023.
Industry bodies argue the ceiling suppresses the merchant case for storage, and
a recent peer-reviewed study identifies it as a constraint, but the magnitude
has not been measured. We measure it. Applying each ceiling to one observed
hourly IEX series and solving the storage dispatch exactly by dynamic
programming, a co-located four-hour battery loses 18.6\% of its arbitrage value
between the Rs20 and Rs10 ceilings, while the capture price of the solar plant
beside it falls 4.3\%. The asymmetry is mechanical: solar output concentrates in
midday hours the ceiling never reaches, while a battery earns in the evening
peak, and the Rs10 ceiling censors 5.7\% of hours in that series while removing
close to a fifth of the arbitrage value. Because a binding ceiling censors rather
than estimates, revenue computed on the capped series is a lower bound on
uncapped revenue and no counterfactual price has to be invented. CERC has since
carved storage out through a high-price segment with a Rs50 ceiling, which
concedes the mechanism and bounds what we measure to the ordinary day-ahead
market. A ceiling defended as consumer protection falls hardest on the
technology that would reduce the scarcity the ceiling exists to suppress.
\end{abstract}

\begin{IEEEkeywords}
electricity markets, price caps, battery storage, arbitrage, solar
co-location, India, techno-economic analysis
\end{IEEEkeywords}

\section{Introduction}

A wholesale price ceiling is a blunt instrument aimed at a specific harm. When
scarcity drives clearing prices to levels that look punitive, a regulator can cap
them, and consumers are spared the tail. India's Central Electricity Regulatory
Commission (CERC) did this twice in quick succession: the ceiling on the power
exchanges fell from Rs20/kWh to Rs12/kWh on 1 April 2022, during the coal supply
crisis, and to Rs10/kWh on 1 April 2023.

That intervention's cost is not borne evenly. A battery earns by buying cheap
and selling dear, which means that the hours generating its revenue are, by
construction, the hours at the top of the price distribution, and a ceiling
placed on that distribution removes them first and removes them entirely.

Indian industry bodies have made this argument, and so, in a peer-reviewed
setting co-authored from within CERC itself, have Suski et al.
\cite{suski2026}, who conclude that storage profitability ``will further
deteriorate if market rules constrain their movement through low price caps or
other `remedial' regulatory measures.'' What no study we are aware of does is
vary the ceiling and measure the value it removes. The mechanism is asserted by
the regulator's own analysis and left unquantified, and it could be wrong in two
ways worth checking before building on it: if the ceilings bind only rarely, or
bind in hours when a battery would not have been discharging anyway, the effect
on the merchant case is negligible.

It is not negligible. We take the hourly Indian Energy Exchange day-ahead series
for 2021 to 2023, apply each ceiling to the same observed prices, and solve the
storage dispatch exactly. A four-hour battery loses \textbf{18.6\%} of its
arbitrage value between the Rs20 and Rs10 ceilings; a two-hour battery loses
\textbf{20.9\%}. Over the same series and the same ceilings, the capture price of
a co-located solar plant falls \textbf{4.3\%}.

That asymmetry is the paper's central result, and it is more robust than either
number alone, because whatever one assumes about how bidding behaviour might
respond to a different ceiling, the assumption applies to both technologies on
the same price series and therefore cancels from the comparison. The mechanism
is a matter of timing. Solar generates at midday, when Indian day-ahead prices
sit near their daily low and no ceiling is within reach; a battery earns after
sunset, in the part of the distribution a ceiling removes first.

Three further results follow, and one boundary on all of them. The ceilings bind
where it costs most: imposing Rs10 on our base series censors only 5.7\% of
hours yet removes 18.6\% of four-hour arbitrage value. For solar, location
matters more than the ceiling on revenue and less on the price received per
unit, so the comparison reverses between the two, while the realised revenue
comparison runs the other way from the counterfactual one, which is precisely
why the fixed-series design is the identifying one. The boundary is narrower. A
high-price segment created in February 2023 lets storage sell at a Rs50 ceiling
(Section~\ref{sec:hpdam}), so what follows measures the ordinary day-ahead
ceiling and not the whole regulatory constraint on merchant storage.

\section{Related Work}

The value of storage arbitrage in a wholesale market is a well-developed
literature. Sioshansi et al. \cite{sioshansi2009} established the standard
approach on PJM data, optimising dispatch under perfect foresight and reporting
the welfare consequences; subsequent work has extended it across US real-time
markets \cite{bradbury2014}, the Australian energy-only market
\cite{mcconnell2015}, and European day-ahead markets at continental scale
\cite{zafirakis2016, brijs2019, mercier2023, komorowska2024}. Methodological
refinements have addressed degradation \cite{wankmuller2017}, dynamic efficiency
and the choice of optimisation formulation \cite{grimaldi2024}, and the limits
that forecast error places on realisable arbitrage \cite{antweiler2021}.

Price caps have their own literature, largely concerned with resource adequacy:
a ceiling set below the value of lost load truncates scarcity rents and produces
the missing-money problem, so capacity mechanisms or scarcity pricing reforms are
needed to restore the investment signal \cite{wilson2000, billimoria2019,
papavasiliou2020, woo2019}, and caps interact with capacity payments and with
renewable-driven volatility \cite{milstein2019, milstein2015}. This work treats
generation adequacy in general. It does not isolate storage, and it does not
measure what a specific ceiling removes from a specific asset.

A third literature concerns the market value of variable renewables and its
decline with penetration \cite{hirth2013, brown2021}, storage's role in that
setting including its market power \cite{williams2022}, co-location design and
planning \cite{park2021}, and the price behaviour of the Indian exchange
\cite{chattopadhyay2023}.

The gap this paper fills sits at the intersection. The storage arbitrage
literature takes the price series as given and does not vary the regulatory
ceiling, while the price cap literature treats adequacy in aggregate. Suski et
al. \cite{suski2026} bring the two together in the Indian setting, and their
treatment of the ceiling is the closest prior work to ours: they document that
prices rested at the cap for more than a tenth of the hours in 2022 and 2023,
and they name low price caps as a threat to the storage business case. Their
object, however, is a comparison of pumped hydro, battery and run-of-river
options, and their sensitivity analysis varies capital cost rather than the
ceiling, which enters as an observed feature of the market and a stated risk
rather than a quantity whose cost is measured. We vary it while holding
everything else fixed.

\section{Institutional Setting}

The Indian Energy Exchange operates a day-ahead market with hourly clearing and a
single national price in the absence of congestion. CERC sets the maximum
clearable price administratively.

During the coal supply crisis of early 2022, with exchange prices repeatedly
touching the then-prevailing Rs20/kWh ceiling, CERC reduced the ceiling to
Rs12/kWh with effect from 1 April 2022. A second reduction to Rs10/kWh followed
on 1 April 2023. Both interventions were framed as protecting distribution
utilities and, through them, consumers from extreme prices during a supply
emergency.

For a merchant storage asset these are not distributional details. The ceiling
determines the maximum spread available on any day in this segment, and for a
merchant battery the spread is the revenue line.

\subsection{The high-price carve-out}
\label{sec:hpdam}

CERC subsequently created an exception, and it matters for how the results below
should be read. By order of 16 February 2023 the Commission approved a High
Price Day Ahead Market within the integrated day-ahead segment, restricted to
high-variable-cost sellers: gas plants running on imported regasified LNG or
naphtha, imported-coal plants, and battery energy storage systems. Its ceiling
was Rs50/kWh against the Rs12/kWh then applying to the ordinary day-ahead
market, and the exchange began trading it that March.

Two things follow. The first is that storage was explicitly carved out, which is
the regulator conceding the mechanism this paper measures: the ordinary ceiling
was recognised as an obstacle for precisely the class of seller whose revenue
comes from the upper tail. The second is a scope condition on our results. From
March 2023 a merchant battery had a legal route to sell above the ordinary
ceiling, so what we measure is the value the \emph{day-ahead} ceiling removes,
not the total constraint on Indian merchant storage. We return to this in
Section~VII.

\section{Data}

\subsection{Prices}

Hourly IEX prices are taken from a published, citable dataset
\cite{veeramsetty2024} released under CC BY 4.0 with a DOI, comprising 23,304
hourly observations. We use a published dataset rather than scraping the exchange
because it is stable, licensed and reproducible by a reviewer.

The file carries no timestamp column, so the series was reconstructed and then
validated rather than assumed. The dataset's own lagged price columns supply the
test, the one-hour lag matching the previous row for 100.00\% of 23,303 rows and
the 24-hour lag for 100.00\% of applicable rows, which establishes a strictly
contiguous hourly series. The start date follows from the row count. Stated
coverage from 1 January 2021 to 31 August 2023 contains 23,352 hours against
23,304 rows, a difference of exactly the 48 leading hours consumed in
constructing the 48-hour lag feature, so the series begins at 3 January 2021.
One day, 2 April 2022, is excluded for straddling the first ceiling change.

Two independent checks confirm these are day-ahead prices and that the
reconstruction is correctly aligned. Aggregating over the Indian fiscal year
2022--23 gives a mean of Rs5.94/kWh, matching the published IEX day-ahead average
for that year exactly; a misaligned reconstruction would shift the fiscal window
and break the match. Second, the series reproduces the CERC interventions
unaided. The maximum price is Rs20,000/MWh until the delivery hour of
1 April 2022, Rs12,000/MWh through to March 2023, and Rs10,000/MWh thereafter.
The ceilings appear at the right values, in the right order, on the right dates.

\subsection{Solar}

Hourly solar output is from PVGIS v5.3 \cite{pvgis}, the European Commission
Joint Research Centre's satellite-derived irradiance service, for 2021 to 2023 at
1 kWp with 14\% system loss. Eight sites were selected to span India's solar
geography: Jaisalmer, Bhadla and Jodhpur in Rajasthan, Bhuj in Gujarat, Pavagada
and Bengaluru in Karnataka, Kurnool in Andhra Pradesh, and Nagpur in
Maharashtra. Mean annual yield ranges from 1,212 to 1,407 kWh/kWp.

\section{Method}

\subsection{Identification}

The obvious objection to any study of a price ceiling is that the counterfactual
uncapped price is unknowable. We do not estimate one.

A binding ceiling is a censoring event. In an hour that clears at the cap, the
true market-clearing price was at least the cap and possibly far above it, which
means that arbitrage revenue computed on a capped series is a \emph{lower bound}
on the revenue the same asset would have earned without the cap. The quantity we
report is therefore bounded from below by construction, and the bound tightens
rather than loosens as the ceiling binds more often.

The bound is conditional. It holds only if offer behaviour is invariant to the
ceiling, and the condition belongs here rather than among the limitations,
because should a lower ceiling change how participants bid, prices in hours that
never touch the cap move as well and censoring no longer isolates the ceiling
from anything else. Our design cannot test this: holding the series fixed is
what buys the bound in the first place. The direction of any such effect is not
obvious either, since a ceiling removing the top of the distribution weakens the
incentive to bid into it at all.

Regime boundaries are assigned from the series rather than from the order dates,
each regime beginning on the first day after which its ceiling is never exceeded.
This places the A/B boundary at 1 April 2022 and the B/C boundary at 30 March
2023, both within a day of CERC's interventions, and yields 10,896, 8,664 and
3,720 hours respectively. Two transition days fall between regimes and are
excluded.

The comparison is run on a single observed series. Regime A, the period before
April 2022, is the least censored series available, with the ceiling binding in
0.44\% of hours. We apply each candidate ceiling to that same series, so the
ceiling is the only quantity that differs between columns and no other market
condition moves with it.

\subsection{Dispatch}

Storage dispatch is solved exactly, by dynamic programming over state of charge,
under price-taking perfect foresight on a daily horizon, with power of 1 MW,
energy of 2 or 4 MWh, round-trip efficiency of 85\%, and the battery beginning
and ending each day empty. Multiple cycles per day are permitted where
profitable. Every schedule reported is chronologically feasible, which we verify
rather than assume: the check passes for 100\% of dispatched days.

Perfect foresight yields an upper bound on realisable arbitrage, which is the
convention in this literature \cite{sioshansi2009, antweiler2021}. Since our
object is a ratio between two ceilings applied to the same series and the same
solver, the foresight assumption affects both sides equally.

The daily horizon was tested against longer ones. Sioshansi et al. optimise two
weeks at a time, which by their account admits both intra-day and inter-day
arbitrage, including heavier charging at weekends. Extending our horizon to a
week or a fortnight raises four-hour revenue by about 2\%, from Rs5.246M to
Rs5.371M per MW-year, and moves the value removed at the Rs10 ceiling from
18.6\% to 19.2\%. The result is not an artifact of horizon choice, and the daily
horizon is the conservative option.

Our round-trip efficiency of 85\% is higher than the 80\% Sioshansi et al.
assume, reflecting lithium-ion rather than the pumped hydro of their benchmark.
A higher efficiency raises the level of arbitrage revenue on both sides of every
comparison we report, and the quantity of interest here is the ratio between
ceilings rather than the level.

\subsection{Solar capture}

Solar capture price is the generation-weighted mean price received, computed
hourly for each site under each ceiling. The capture ratio is that price divided
by the simple time-weighted mean.

\section{Results}
\label{sec:results}

\begin{figure*}[t]
\centering
\includegraphics[width=0.85\textwidth]{figB1.pdf}
\caption{Panel (a): price duration curves by ceiling regime. The flat upper
segments under the Rs12 and Rs10 ceilings are the censored hours. Panel (b):
share of hours clearing at the ceiling, which rises from 0.44\% to 19.2\% as the
ceiling falls.}
\label{fig:censoring}
\end{figure*}

\subsection{The ceilings bind, and they bind in the tail}

Table~\ref{tab:regimes} gives the three regimes as they appear in the data.
Under the Rs20 ceiling the cap is nearly irrelevant, binding in 0.44\% of hours.
Under Rs12 it binds in 12.4\%, and under Rs10 in 19.2\%, so close to one hour in
five clears at the administrative maximum.

\begin{table}[htbp]
\caption{Price regimes and censoring}
\label{tab:regimes}
\centering
\begin{tabular}{lrrrr}
\hline
\textbf{Regime} & \textbf{Ceiling} & \textbf{Hours} & \textbf{Mean} & \textbf{At cap}\\
\hline
A & Rs20/kWh & 10,896 & Rs4.26 & 0.44\%\\
B & Rs12/kWh & 8,664 & Rs5.92 & 12.40\%\\
C & Rs10/kWh & 3,720 & Rs5.39 & 19.19\%\\
\hline
\end{tabular}
\end{table}

In both tighter regimes the 99th percentile of the price distribution equals the
ceiling exactly, against Rs18.72 under the Rs20 ceiling. This follows
mechanically once the censored share passes one percent and is reported as a
description of the shape rather than as independent evidence: above the
99th percentile there is no variation left to price, and the censoring falls on
exactly the hours a battery exists to sell into.

\subsection{Value removed from storage}

Table~\ref{tab:removed} reports arbitrage revenue under each ceiling applied to
the regime A series. A four-hour battery earns Rs5.246M per MW-year under the
Rs20 ceiling and Rs4.273M under Rs10, which is a loss of 18.6\%, while a
two-hour battery gives up 20.9\% across the same range. The shorter duration
loses proportionally more. Its revenue is concentrated in fewer and
higher-priced hours, so truncating the top of the distribution takes a larger
share of it.

\begin{table}[htbp]
\caption{Arbitrage value removed by the ceiling, regime A series}
\label{tab:removed}
\centering
\begin{tabular}{lrrrrr}
\hline
& \textbf{Hours} & \multicolumn{2}{c}{\textbf{2-hour}} & \multicolumn{2}{c}{\textbf{4-hour}}\\
\textbf{Ceiling} & censored & Rs M & removed & Rs M & removed\\
\hline
Rs20 & 0.44\% & 3.583 & 0.0\% & 5.246 & 0.0\%\\
Rs12 & 4.24\% & 3.069 & 14.4\% & 4.588 & 12.5\%\\
Rs10 & 5.69\% & 2.835 & 20.9\% & 4.273 & 18.6\%\\
Rs8 & 8.03\% & 2.531 & 29.4\% & 3.847 & 26.7\%\\
\hline
\end{tabular}
\end{table}

The censored column is important for reading the rest of the table correctly.
Applying the Rs10 ceiling to the regime A series censors 5.69\% of hours, not the
19.2\% observed under the realised Rs10 regime in Table~\ref{tab:regimes}. These
are different quantities. One is the effect of a ceiling on a given price
series; the other is what happened when a tighter ceiling met a market whose
underlying volatility had risen, and a battery loses close to a fifth of its
arbitrage value to a ceiling binding in about one hour in eighteen because the
hours it binds in are the ones that matter most.

We include a hypothetical Rs8 ceiling because the direction of Indian policy over
the sample period was downward, and because the relationship is visibly convex,
in Table~\ref{tab:removed}: each successive reduction removes more value
than the one before it.


\subsection{The asymmetry with solar}

Under the same series and the same ceilings, solar capture price falls from
Rs4.89/kWh to Rs4.68/kWh, a decline of 4.3\%. The capture ratio sits between 0.95
and 0.98 throughout, meaning Indian merchant solar in this window received
close to the flat time-weighted average price.

Between the same two ceilings, then, a four-hour battery gives up roughly four
times what the solar plant beside it gives up. No modelling is required to see why: solar output is concentrated in the middle
of the day, Indian day-ahead prices are at their lowest in exactly those hours,
and the ceiling never comes close to binding on them, whereas a battery earns
its revenue after sunset, in the evening peak, which is the part of the
distribution the ceiling was designed to remove.

Across the eight sites, mean capture price varies by 2.3\%, against the 4.3\%
the ceiling moves it. That comparison does not survive restatement as revenue,
since it holds output fixed while annual yield ranges from 1,212 to 1,407
kWh/kWp across the same sites, and revenue is yield multiplied by capture price.
On revenue the ordering reverses: moving between these sites is worth 12.6\%
against the ceiling's 4.3\%, so location dominates by roughly three to one, and
the useful statement is therefore the narrow one, that the ceiling is the larger
influence on what a solar plant is paid per unit sold while location remains the
larger influence on what it earns. For storage the distinction does not arise.
Arbitrage value depends on the shape of the price distribution rather than on
any generation profile of its own.

\subsection{Realised revenue runs the other way}

The counterfactual comparison and the realised one point in opposite directions,
and we report this because it is the most obvious way the result could be
misread.

Measured on the prices that actually obtained in each regime, a four-hour
battery earned Rs5.246M per MW-year in regime A, Rs7.572M in regime B and
Rs6.356M in regime C. Realised revenue was \emph{higher} under the tighter
ceilings. The explanation is the coal supply crisis, since through 2022 the
level and volatility of Indian day-ahead prices rose considerably further than
the ceiling removed, and restricting to the matched months of April to August
preserves that ordering.

This does not contradict the counterfactual result. It is why the counterfactual
design is necessary, since comparing regimes as they occurred confounds the
ceiling with everything else that changed in the Indian power system between
2021 and 2023, whereas holding the price series fixed and varying only the
ceiling isolates it. The realised series misleads. A reader taking it at face
value would conclude that tightening the cap raised storage revenue, an artifact
of the fuel crisis and not a property of price ceilings.

\subsection{Degradation and forecast error}

Two assumptions in the dispatch carry real weight, and both can be measured
rather than argued.

The first concerns profit. We report gross arbitrage revenue, whereas a battery
that cycles pays for the cycling. Rather than importing a degradation cost from
another market, we sweep one through the optimisation, charging it against every
MWh discharged. In rupees the ceiling becomes cheaper as cycling becomes dearer,
taking Rs0.974M per MW-year at zero cost against Rs0.852M at Rs2,000/MWh,
because an expensive cycle is one the battery declines to run: daily cycling
falls from 6.09 to 5.07 under the Rs20 ceiling across that range. As a
proportion, which is how we report it, the quantity moves the opposite way.
Table~\ref{tab:deg} gives the sweep. At Rs500/MWh---roughly an eighth of this
series' mean price---the Rs10 ceiling takes 22.3\% of four-hour arbitrage value
against 18.6\% with no cycling charge, and by Rs2,000/MWh it takes 38.7\%. The
compressions multiply. A ceiling narrows the spread a battery earns upon, and a
fixed charge against throughput then consumes a larger share of whatever
survives, which leaves our headline proportion conservative with respect to
degradation rather than generous.

\begin{table}[htbp]
\caption{Arbitrage value net of a marginal cycling cost, four-hour battery,
regime A series, Rs M per MW-year}
\label{tab:deg}
\centering
\begin{tabular}{lrrr}
\hline
\textbf{Cost, Rs/MWh discharged} & \textbf{Rs20} & \textbf{Rs10} & \textbf{Removed}\\
\hline
0 & 5.246 & 4.273 & 18.6\%\\
500 & 4.232 & 3.289 & 22.3\%\\
1,000 & 3.396 & 2.481 & 26.9\%\\
2,000 & 2.201 & 1.349 & 38.7\%\\
\hline
\end{tabular}
\end{table}

The second is perfect foresight. We re-solve the dispatch against a forecast an
operator could genuinely have held, taking the previous day's realised prices as
the curve, committing to the resulting schedule, and settling it at whatever
prices then obtained, with the same ceiling imposed on forecast and settlement
alike. Across the 453 days carrying a predecessor that operator captures 89.6\%
of perfect-foresight revenue with a four-hour battery and 87.6\% with a two-hour
one. The ratio barely moves. What the Rs10 ceiling withdraws goes from 18.6\% to
17.6\% for the four-hour battery and from 20.9\% to 19.1\% for the two-hour, so
forecast error costs roughly a tenth of the level against a single point of the
ratio. That ordering is what the design predicts, since both scenarios share a
forecast of identical quality and the error very largely divides out.

\section{Discussion}

Instrument and target are misaligned here in a way this measurement makes
specific. A price ceiling exists to suppress scarcity prices. Storage reduces
scarcity directly, by moving energy out of surplus hours and into the hours when
the system is short, which is the same thing as flattening the price
distribution the ceiling is there to truncate. The ceiling removes 18.6\% of the
revenue of the asset that would do this, and 4.3\% of the revenue of the asset
that would not. The intervention falls hardest on the one technology---the
flexible one---whose deployment would make the intervention less necessary.

This is not an argument that the ceilings were wrong. A regulator facing a
supply emergency with distribution utilities in financial distress works on a
shorter timescale than storage investment. It is an argument that the
instrument's cost has a specific incidence, unquantified until now, and that
CERC's own high-price carve-out acknowledges exactly that.

Magnitudes of this size matter to an investment case. A four-hour merchant
battery in India surrenders close to a fifth of its arbitrage revenue to a
parameter that is set administratively, was revised twice inside two years, and
sits on the revenue line rather than the cost line, where it is correspondingly
difficult to hedge and to finance against. The convexity compounds it, since the
move from Rs12 to Rs10 removed more value than the larger move from Rs20 to
Rs12. Any jurisdiction weighing a ceiling against a storage deployment target is
choosing between them to a degree this measurement makes quantitative, and the
generalisable instruction is to price that trade rather than to assume the two
policies are independent.

\subsection{Limitations}

We report gross arbitrage revenue rather than profit, and Section~VI-E measures
what cycling cost does to that \cite{wankmuller2017}. The rupees removed fall as
cycling grows dearer while the share removed rises, so the percentages reported
here are a lower bound on the proportional profit impact, as they are again for
the censoring reason given in Section~V-A. We also model the day-ahead segment
alone. A real Indian battery earns across day-ahead, real-time, ancillary
services and contracted capacity, so the share of total revenue exposed to this
particular ceiling is smaller than the share of day-ahead revenue, and two
further simplifications work in the same conservative direction, since a market
clearing at the ceiling is rationed and a battery bidding there may not be fully
dispatched, whereas our solver assumes it is.

The sample runs to August 2023 and covers one exchange and one market design.
The lower-bound argument establishes that the value removed is at least what we
report, but not how much more it might be, since that would require knowing how
high uncapped prices would have gone. A deeper limit sits behind that one.
Bidding behaviour under a different ceiling may itself differ, which our design
holds fixed by construction and cannot test, and lifting that restriction needs
a structural model of offer behaviour rather than additional data. Two
extensions are straightforward by comparison: participation across real-time,
ancillary and contracted segments alters a battery's exposure to any single
ceiling, and the Rs50 segment has accumulated enough history to be measured
rather than described.

\section{Conclusion}

India's day-ahead price ceiling fell from Rs20/kWh to Rs10/kWh across two
interventions inside our sample. Applying each ceiling to a single observed price
series and solving storage dispatch exactly, the ceiling removes 18.6\% of the
arbitrage value of a four-hour battery and 4.3\% of the capture value of the
solar plant co-located with it. Neither figure is an artifact of the modelling
choices that produce it: a marginal cycling cost raises the share removed rather
than lowering it, and an operator dispatching against the previous day's prices
loses about a tenth of the level while the ratio between ceilings moves by a
single point.

The claim that price ceilings constrain the merchant storage case has been made
repeatedly by industry and identified in the literature. It is correct, and the
magnitude is now measured: close to a fifth of arbitrage revenue between the
ceilings India actually set, concentrated on the technology whose deployment
would reduce the scarcity that motivated the ceiling originally.

\section*{Data and Reproducibility}

The price dataset \cite{veeramsetty2024} is public under CC BY 4.0 and the solar
service \cite{pvgis} is open. Dispatch and analysis scripts, with the artifacts
behind every figure, are at
\url{https://github.com/jerrysilvon/icisas2026-replication}.

\begin{thebibliography}{00}

\bibitem{sioshansi2009} R. Sioshansi, P. Denholm, T. Jenkin, and J. Weiss, ``Estimating the value of electricity storage in PJM: Arbitrage and some welfare effects,'' \textit{Energy Econ.}, vol. 31, no. 2, pp. 269--277, 2009.

\bibitem{bradbury2014} K. Bradbury, L. Pratson, and D. Pati\~no-Echeverri, ``Economic viability of energy storage systems based on price arbitrage potential in real-time U.S. electricity markets,'' \textit{Appl. Energy}, vol. 114, pp. 512--519, 2014.

\bibitem{mcconnell2015} D. McConnell, T. Forcey, and M. Sandiford, ``Estimating the value of electricity storage in an energy-only wholesale market,'' \textit{Appl. Energy}, vol. 159, pp. 422--432, 2015.

\bibitem{zafirakis2016} D. Zafirakis, K. J. Chalvatzis, G. Baiocchi, and G. Daskalakis, ``The value of arbitrage for energy storage: Evidence from European electricity markets,'' \textit{Appl. Energy}, vol. 184, pp. 971--986, 2016.

\bibitem{brijs2019} T. Brijs, F. Geth, C. De Jonghe, and R. Belmans, ``Quantifying electricity storage arbitrage opportunities in short-term electricity markets in the CWE region,'' \textit{J. Energy Storage}, vol. 25, p. 100899, 2019.

\bibitem{mercier2023} T. Mercier, M. Olivier, and E. De Jaeger, ``The value of electricity storage arbitrage on day-ahead markets across Europe,'' \textit{Energy Econ.}, vol. 123, p. 106721, 2023.

\bibitem{komorowska2024} A. Komorowska and P. Olczak, ``Economic viability of Li-ion batteries based on the price arbitrage in the European day-ahead markets,'' \textit{Energy}, vol. 290, p. 130009, 2024.

\bibitem{wankmuller2017} F. Wankm\"uller, P. R. Thimmapuram, K. G. Gallagher, and A. Botterud, ``Impact of battery degradation on energy arbitrage revenue of grid-level energy storage,'' \textit{J. Energy Storage}, vol. 10, pp. 56--66, 2017.

\bibitem{grimaldi2024} A. Grimaldi, F. D. Minuto, J. Brouwer, and A. Lanzini, ``Profitability of energy arbitrage net profit for grid-scale battery energy storage considering dynamic efficiency and degradation,'' \textit{J. Energy Storage}, vol. 95, p. 112380, 2024.

\bibitem{antweiler2021} W. Antweiler, ``Microeconomic models of electricity storage: Price forecasting, arbitrage limits, curtailment insurance, and transmission line utilization,'' \textit{Energy Econ.}, vol. 101, p. 105390, 2021.

\bibitem{wilson2000} J. F. Wilson, ``Scarcity, market power, and price caps in wholesale electric power markets,'' \textit{Electr. J.}, vol. 13, no. 9, pp. 33--46, 2000.

\bibitem{billimoria2019} F. Billimoria and R. Poudineh, ``Market design for resource adequacy: A reliability insurance overlay on energy-only electricity markets,'' \textit{Util. Policy}, vol. 60, p. 100935, 2019.

\bibitem{papavasiliou2020} A. Papavasiliou, ``Scarcity pricing and the missing European market for real-time reserve capacity,'' \textit{Electr. J.}, vol. 33, no. 10, p. 106863, 2020.

\bibitem{woo2019} C. K. Woo, I. Milstein, A. Tishler, and J. Zarnikau, ``A wholesale electricity market design sans missing money and price manipulation,'' \textit{Energy Policy}, vol. 134, p. 110988, 2019.

\bibitem{milstein2019} I. Milstein and A. Tishler, ``On the effects of capacity payments in competitive electricity markets: Capacity adequacy, price cap, and reliability,'' \textit{Energy Policy}, vol. 129, pp. 370--385, 2019.

\bibitem{milstein2015} I. Milstein and A. Tishler, ``Can price volatility enhance market power? The case of renewable technologies in competitive electricity markets,'' \textit{Resour. Energy Econ.}, vol. 41, pp. 70--90, 2015.

\bibitem{hirth2013} L. Hirth, ``The market value of variable renewables,'' \textit{Energy Econ.}, vol. 38, pp. 218--236, 2013.

\bibitem{brown2021} T. Brown and L. Reichenberg, ``Decreasing market value of variable renewables can be avoided by policy action,'' \textit{Energy Econ.}, vol. 100, p. 105354, 2021.

\bibitem{williams2022} O. Williams and R. Green, ``Electricity storage and market power,'' \textit{Energy Policy}, vol. 164, p. 112872, 2022.

\bibitem{park2021} H. Park, ``A stochastic planning model for battery energy storage systems coupled with utility-scale solar photovoltaics,'' \textit{Energies}, vol. 14, no. 5, p. 1244, 2021.

\bibitem{chattopadhyay2023} D. Chattopadhyay, S. K. Chatterjee, and S. K. Soonee, ``Spotlight on the spot market: A review of the Indian wholesale electricity market,'' \textit{Electr. J.}, vol. 36, no. 1, p. 107239, 2023.

\bibitem{suski2026} A. Suski, I. Deleque Curiel, D. Chattopadhyay, S. Chatterjee, and D. Malovic, ``Evaluating storage profitability: A market-based comparison of hydro and battery options in India,'' \textit{J. Energy Storage}, vol. 144, p. 119803, 2026.

\bibitem{veeramsetty2024} V. Veeramsetty, S. R. Salkuti, B. Nikitha, and T. Siddartha, ``Electrical energy price dataset: IEX data,'' Mendeley Data, V1, 2024, doi: 10.17632/v5znbkzjd4.1.

\bibitem{pvgis} European Commission Joint Research Centre, ``Photovoltaic Geographical Information System (PVGIS), v5.3,'' 2024.

\end{thebibliography}

\end{document}
